% AirESP32ace
% Created by VM, August 2026

% Build with XeLaTeX on macOS. Avenir Next is supplied by macOS and is not
% redistributed with the project.
\documentclass[10pt,letterpaper]{article}

\usepackage[margin=0.62in]{geometry}
\usepackage{fontspec}
\usepackage{graphicx}
\usepackage[table]{xcolor}
\usepackage{tikz}
\usepackage{tabularx}
\usepackage{booktabs}
\usepackage{enumitem}
\usepackage{listings}
\usepackage{microtype}
\usepackage{hyperref}
\usepackage{ragged2e}

\setmainfont{Avenir Next}
\setsansfont{Avenir Next}
\setmonofont{Menlo}
\renewcommand{\familydefault}{\sfdefault}

\definecolor{Canvas}{HTML}{080A0C}
\definecolor{Panel}{HTML}{111519}
\definecolor{PanelTwo}{HTML}{171C20}
\definecolor{Text}{HTML}{F4F7F5}
\definecolor{Muted}{HTML}{98A19C}
\definecolor{Line}{HTML}{29312D}
\definecolor{Accent}{HTML}{20D45B}
\definecolor{AccentDark}{HTML}{0D6D35}
\definecolor{CodeBG}{HTML}{0C100E}

\pagecolor{Canvas}
\color{Text}
\hypersetup{
  colorlinks=true,
  urlcolor=Accent,
  linkcolor=Accent,
  pdftitle={AirESP32ace - Getting Started},
  pdfauthor={AirESP32ace}
}
\setlength{\parindent}{0pt}
\setlength{\parskip}{4pt}
\pagestyle{empty}
\emergencystretch=2em

\lstdefinestyle{terminal}{
  basicstyle=\ttfamily\fontsize{8.2}{10.2}\selectfont\color{Text},
  backgroundcolor=\color{CodeBG},
  frame=single,
  rulecolor=\color{Line},
  framerule=0.5pt,
  framesep=7pt,
  xleftmargin=0pt,
  xrightmargin=0pt,
  aboveskip=5pt,
  belowskip=5pt,
  showstringspaces=false,
  columns=fullflexible,
  keepspaces=true
}

\newcommand{\Kicker}[1]{%
  {\fontsize{7.6}{9}\selectfont\bfseries\color{Accent}\MakeUppercase{#1}\par}%
}
\newcommand{\SectionTitle}[1]{%
  {\fontsize{14.5}{17}\selectfont\bfseries\color{Text}#1\par}%
  \vspace{2pt}\textcolor{Line}{\rule{\linewidth}{0.55pt}}\vspace{4pt}%
}
\newcommand{\SmallTitle}[1]{%
  {\fontsize{9.3}{11}\selectfont\bfseries\color{Text}#1\par}%
}
\newcommand{\Footer}[1]{%
  \vfill
  \textcolor{Line}{\rule{\linewidth}{0.5pt}}\vspace{4pt}
  {\fontsize{7.2}{8.6}\selectfont\color{Muted}
    AirESP32ace \hfill Getting Started \hfill #1 / 3\par}
}
\newcommand{\Service}[3]{%
  \begin{minipage}[t]{0.31\linewidth}
    \parbox[c][0.76in][c]{\linewidth}{%
      \Kicker{#1}
      {\fontsize{8.1}{10.1}\selectfont\color{Muted}#2\par}
      \vspace{2pt}{\fontsize{7.1}{8.6}\selectfont\color{Accent}#3\par}%
    }
  \end{minipage}%
}
\newcommand{\Step}[3]{%
  \begin{minipage}[t]{\linewidth}
    \begin{tabularx}{\linewidth}{@{}p{0.36in}X@{}}
      \raisebox{-0.12in}{\tikz\node[circle,fill=Accent,text=Canvas,
        font=\bfseries\fontsize{9}{10}\selectfont,minimum size=0.27in,
        inner sep=0pt]{#1};} &
      {\fontsize{9.2}{11}\selectfont\bfseries\color{Text}#2}\par
      {\fontsize{7.8}{9.7}\selectfont\color{Muted}#3}\tabularnewline
    \end{tabularx}
  \end{minipage}\vspace{4pt}%
}
\newcommand{\GuideShot}[4]{%
  \begin{minipage}[t]{0.485\linewidth}
    \centering
    \includegraphics[width=#4]{#1}\par
    \vspace{3pt}
    \raggedright
    {\fontsize{8.4}{10.2}\selectfont\bfseries\color{Text}#2\par}
    {\fontsize{7.2}{8.8}\selectfont\color{Muted}#3\par}
  \end{minipage}%
}

\begin{document}

% ---------------------------------------------------------------------------
% Page 1
% ---------------------------------------------------------------------------
\begin{minipage}[t]{0.74\linewidth}
  \vspace{0pt}
  \Kicker{Open source / quick start}
  \vspace{5pt}
  {\fontsize{29}{31}\selectfont\bfseries AirESP32ace\par}
  \vspace{3pt}
  {\fontsize{13}{16}\selectfont\color{Muted}From clone to the first live screen\par}
  \vspace{7pt}
  {\fontsize{8.6}{11}\selectfont\color{Text}
    ESP32-S3 firmware and a native macOS companion for the Waveshare
    1.75-inch circular AMOLED board. This guide covers the shortest reliable
    path to a working device.\par}
\end{minipage}\hfill
\begin{minipage}[t]{0.20\linewidth}
  \vspace{-4pt}
  \includegraphics[width=\linewidth]{assets/AirESP32ace_AppIcon_Transparent.png}
\end{minipage}

\vspace{7pt}
\textcolor{Accent}{\rule{\linewidth}{1.4pt}}
\vspace{9pt}

\begin{minipage}[t]{0.47\linewidth}
  \SectionTitle{What you need}
  \begin{itemize}[leftmargin=0.18in,itemsep=2.5pt,topsep=1pt]
    \item \textbf{Waveshare ESP32-S3-Touch-AMOLED-1.75C} with 32 MB flash and 8 MB PSRAM.
    \item A \textbf{data-capable USB-C cable}; charging-only cables cannot flash the board.
    \item A Mac on \textbf{macOS 15 or newer}, Git, and Xcode Command Line Tools.
    \item 2.4 GHz Wi-Fi. An iPhone hotspot can be added as an optional second network.
    \item Optional provider accounts for richer aircraft and tracked-flight data.
  \end{itemize}
  \vspace{3pt}
  \SmallTitle{No Apple Developer account required}
  {\fontsize{7.6}{9.5}\selectfont\color{Muted}
    Development builds use local ad-hoc signing. API credentials stay on
    your Mac or device and are never included in the repository.\par}
\end{minipage}\hfill
\begin{minipage}[t]{0.47\linewidth}
  \SectionTitle{What runs where}
  \begin{tabularx}{\linewidth}{@{}p{0.76in}X@{}}
    \textcolor{Accent}{\bfseries Device} &
      ESP-IDF 5.5+, LVGL 9.5, the Waveshare board component, live UI,
      OpenSky radar, weather, ISS, and cached RGB565 artwork.\\[5pt]
    \textcolor{Accent}{\bfseries Mac} &
      Native SwiftUI companion, MapKit renderers, setup and firmware tools,
      tracked-flight scheduling, API validation, and the exact LVGL digital twin.\\[5pt]
    \textcolor{Accent}{\bfseries Sync} &
      Selected cities, settings, flight plans, maps, and generated caches are
      sent to one paired device. No cloud account is required.\\
  \end{tabularx}
  \vspace{4pt}
  {\fontsize{7.4}{9.3}\selectfont\color{Muted}
    The firmware deliberately separates fast animation from slower network
    updates and uses PSRAM, cached assets, and partial redraws to keep the
    466 by 466 AMOLED responsive.\par}
\end{minipage}

\vspace{10pt}
\SectionTitle{Aircraft data at a glance}
\noindent
\Service{OpenSky}{Nearby aircraft state vectors and best-effort live telemetry.}{Anonymous works; credentials improve limits.}
\hfill
\Service{FlightAware}{Primary provider for a selected scheduled flight and its status.}{Optional AeroAPI key; Mac only.}
\hfill
\Service{AirLabs}{Fallback route, metadata, telemetry, and airline identity.}{Optional key; complements OpenSky.}

\vspace{7pt}
{\fontsize{7.4}{9.2}\selectfont\color{Muted}
  Other public services: Open-Meteo for weather and air quality, RainViewer for
  radar frames, Where the ISS at for live ISS state, and Launch Library 2 for
  NASA, Roscosmos, and SpaceX launch schedules. Provider coverage, quotas, and
  terms remain under each provider's control.\par}

\Footer{1}

\newpage

% ---------------------------------------------------------------------------
% Page 2 -- visual setup and first sync. The screen map is appended as page 3
% by the build command so it retains its readable landscape page size.
% ---------------------------------------------------------------------------
\Kicker{First run / macOS companion}
\vspace{3pt}
{\fontsize{22}{25}\selectfont\bfseries From USB-C to the first live screen\par}
{\fontsize{8.5}{10.5}\selectfont\color{Muted}
  The companion handles first-time flashing, Wi-Fi and API configuration, map
  generation, and later Sync operations. Use a data-capable USB-C cable.\par}
\vspace{4pt}

\begin{lstlisting}[style=terminal]
git clone https://github.com/vmalis/AirESP32ace.git
cd AirESP32ace/companion/AirspaceCompanion
swift run AirESP32ace
\end{lstlisting}

\vspace{3pt}
\GuideShot{../../output/guide/companion-screens/device-chooser.png}
  {1. Open the device control}
  {Choose \textbf{Set Up New Device}. Firmware updates later use the adjacent action without replacing settings.}
  {2.15in}
\hfill
\GuideShot{../../output/guide/companion-screens/device-setupConnect.png}
  {2. Connect the tracker}
  {Wait for the actual serial target to appear. First setup and wired recovery require USB-C.}
  {2.15in}

\vspace{7pt}
\GuideShot{../../output/guide/companion-screens/device-setupForm.png}
  {3. Enter configuration}
  {Home Wi-Fi is required for live data. Hotspot and provider credentials are optional.}
  {2.15in}
\hfill
\GuideShot{../../output/guide/companion-screens/device-firmwareSuccess.png}
  {4. Flash and let the board restart}
  {Keep the cable attached through completion. The assistant preserves the submitted network, API, city, and display settings.}
  {2.15in}

\vspace{7pt}
\begin{minipage}{\linewidth}
  \Kicker{First sync}
  {\fontsize{7.8}{9.6}\selectfont\color{Text}
    Wait for \textbf{Reachable} below the digital twin, then choose radar
    cities, clock order, units, Dock views, and up to four non-conflicting
    tracked flights. Press \textbf{Sync}. The Mac generates maps and sends
    configuration plus assets over USB-C when attached, otherwise over Wi-Fi.
    Let the progress bar finish before disconnecting.\par}
  \vspace{3pt}
  {\fontsize{7.2}{8.8}\selectfont\color{Muted}
    No serial target: change the cable or USB port. No live data: verify 2.4 GHz
    Wi-Fi, location, API status, and local provider coverage.\par}%
\end{minipage}

\vspace{6pt}
{\fontsize{7.3}{9.1}\selectfont
  \href{https://github.com/vmalis/AirESP32ace}{Repository}\quad\textcolor{Line}{|}\quad
  \href{https://github.com/vmalis/AirESP32ace/blob/main/docs/guides/AirESP32ace_Device_User_Guide.pdf}{Device guide}\quad\textcolor{Line}{|}\quad
  \href{https://github.com/vmalis/AirESP32ace/blob/main/docs/guides/AirESP32ace_Companion_User_Guide.pdf}{Companion guide}\quad\textcolor{Line}{|}\quad
  \href{https://docs.waveshare.com/ESP32-S3-Touch-AMOLED-1.75C}{Board documentation}\par}

\Footer{2}

\end{document}
